%%% =====================================================================
%%% figset_D.tex  --  TikZ figures, set D of the QCD factorization monograph
%%%   (FIGURES_INDEX_v2 §2: Ch20 heavy, Ch21 power, Ch22 lamet, Ch25 summary;
%%%    §4.2: Ch22 lamet; §4.3: Ch23 yzliu; required and optional figures).
%%% (figures/test_fd2; figure agent figD; local style prefix "fd2")
%%%
%%% Contents (20 figures; in the file they follow the section order of their
%%%   owner chapters; owner chapter in brackets)
%%%   required: fig:heavy-bdpi, fig:heavy-qcdf [Ch20]; fig:power-borel-plane [Ch21];
%%%     fig:lamet-ioffe, fig:lamet-lattice-geometries, fig:lamet-oneloop,
%%%     fig:lattice-matching, fig:lamet-hybrid [Ch22];
%%%     fig:yzliu-soft-hqet, fig:yzliu-theta-strip, fig:yzliu-complex-lorentz,
%%%     fig:yzliu-leading-region, fig:yzliu-threshold [Ch23];
%%%     fig:summary-status-map, fig:summary-hypotheses [Ch25];
%%%   optional: fig:lamet-dependency, fig:lamet-fourier, fig:lamet-validity [Ch22];
%%%     fig:yzliu-single-line, fig:yzliu-borel [Ch23].
%%% The two Ch25 figures are generated from the data of spec_v2/FINAL_OUTLINE_v2
%%%   (master-table rows and tags; incidence data of tab:summary-assumptions,
%%%   bottleneck column) by figures/test_fd2/gen/gen_status_map.py and
%%%   gen_hypotheses.py; rerun them if the outline changes, and check them
%%%   against the generated App P data at integration.
%%%
%%% Requirements: common/preamble_fact.tex (ctexbook, XeLaTeX), which loads tikz
%%%   with decorations.pathmorphing, decorations.markings, arrows.meta, calc,
%%%   positioning, shapes.geometric, patterns, backgrounds, fit, and pgfplots
%%%   (compat=1.17), and contains the style blocks STYLES / STYLES-EXT; plus the
%%%   preamble macros \ii \dd \ee \tp \MSbar \cref \tagR ... \tagV.  No package or
%%%   library is loaded here.  Every extra style is declared in the options of
%%%   its tikzpicture (hence local to that picture) and carries the prefix "fd2".
%%% Each figure is delimited by  %%% BEGIN <label>  /  %%% END <label>.
%%%
%%% Drawing conventions used throughout this set
%%%   * solid line with arrow = quark; the arrow is the fermion-number flow
%%%     (for an antiquark it points against the momentum); thick solid line =
%%%     heavy quark / static colour source;
%%%   * thin separate arrow with label (styles mom/mom') = momentum flow;
%%%   * double line with arrow = Wilson line; the arrow is the path orientation
%%%     of W[b,a] (from a to b), i.e. the colour flow;
%%%   * blob colours as in STYLES-EXT: grey = hard, blue = collinear (first
%%%     direction), orange = collinear (second direction), green = soft,
%%%     light grey = hadron / non-perturbative;
%%%   * tag colours of the rigour levels as in the preamble macros \tagR ...:
%%%     OliveGreen [R], RoyalBlue [A], BurntOrange [P], Plum [C], BrickRed [V].
%%% Cross references to statements, tables and sections of the chapters are
%%% made with \cref (resolved in the book; dummy targets in the test document).
%%% =====================================================================
